\scanpage{120}
% line head | text | 69
\mspage{69}
% line 1 | text | avec
avec
% line 2 | display | F lower shriek between the two Hom categories
\[
F_!:\underline{\operatorname{Hom}}(U^{\mathrm o}\times T,\mathrm{Ens})\longrightarrow
\underline{\operatorname{Hom}}({U'}^{\mathrm o}\times T',\mathrm{Ens}).
\]

% line 3 | text | Posons
\deleted{L'objet} Posons
% line 4 | display | H prime equals F shriek of H
\[
H'=F_!(H):{U'}^{\mathrm o}\times T'\to\mathrm{Ens},
\]
% line 5–7 | text | c'est un bifoncteur bien déterminé
c'est un bifoncteur bien déterminé \deleted{\unsure{étant}} \unsure{la}
seule donnée du diagramme initial \inserted{$(\unsure{*|\alpha})$ p. 67}, ou
ce qui revient au $=$, en termes de
% line 8 | display | U,T,H,f_U,f_T
\[
(U,T,H,f_U:U\unsure{-|\to}U',f_T:T\to T').
\]
% line 9–11 | text | Et on voit que
Et on voit que la donnée d'une flèche
$X^{\unsure{\prime|\prime\prime}}\longrightarrow X''$, $X''$ objet variable de $\mathrm{Cat}$, équivaut
à celle de
% line 12 | display | the two g's and u
\[
g_{U'}:U'-X'',\quad g_{T'}:T'\to X'',\quad
u:H'\to H''=H_{X''}(g_{U'}^{\mathrm o}\times \unsure{g_{T'}|g_{T'}^{\mathrm o}|f_{T'}}),
\]
% line 13–14 | text | c'est à dire, en vertu de 1.3
c'est à dire, en vertu de 1.3, à celle d'une
flèche
% line 14 right | display | X prime subscript nought
\[
X'_0\longrightarrow X'',
\]
% line 15–19 | text | où X prime nought désigne
où $X'_0$ désigne l'objet de $\mathrm{Cat}$ déduit
de \unsure{(}$U',T'$ par recollement au moyen de $H'$.
Cela montre donc que la limite \unsure{inductive}
cherchée n'est autre que $X'_0$. On a
ainsi trouvé le

% line 20–23 | text | Théorème 1. La limite inductive
\uline{Théorème 1}\quad La \deleted{\unsure{somme amalgamée}}\inserted{limite inductive} \inserted{$X'$} dans $\mathrm{Cat}$ du
diagramme $(\alpha)$ page 67, défini en termes
d'un triple $(X,U,T)$ de $\unsure{\mathscr T|\mathcal T}$ et de flèches
$f_U:U-U'$, $f_T:T\to T'$ dans $\mathrm{Cat}$, existe (\unsure{sans mérite|cas mérite|cas mixte})
% archive sheet number: 120
